\chapter{Document Type System}
\label{ch:doctype-system}

OmniLaTeX provides a document type system that maps user-friendly names to
KOMA-Script base classes.  Over 80 aliases resolve to 26 canonical
doctypes, each backed by one of three KOMA classes
(\texttt{scrartcl}, \texttt{scrreprt}, \texttt{scrbook}).

\section{Architecture}

The doctype system uses a three-tier resolution pipeline:

\begin{enumerate}
    \item \textbf{User alias} -- The string passed to
          \texttt{doctype=} (e.g.\ \texttt{paper}, \texttt{talk},
          \texttt{tech-report}).
    \item \textbf{Canonical name} -- The internal doctype identifier
          (e.g.\ \texttt{article}, \texttt{presentation},
          \texttt{technicalreport}).
    \item \textbf{KOMA class} -- The underlying KOMA-Script base class
          (\texttt{scrartcl}, \texttt{scrreprt}, or \texttt{scrbook})
          with associated default options.
\end{enumerate}

This design lets users write intuitive names while the class handles the
complexity of KOMA-Script configuration.  For example,
\texttt{doctype=thesis} resolves through \texttt{thesis} $\to$
\texttt{scrbook} with options
\texttt{listof=totoc,bibliography=totoc,chapterprefix=true,open=right}.

Resolution is implemented via \texttt{\textbackslash omnilatex@matchdoctypes}
in \texttt{omnilatex.cls}, which iterates comma-separated alias lists and
calls \texttt{\textbackslash omnilatex@setdoctype} on the first match.
The doctype string is normalised to lowercase before matching.
Unrecognised doctypes trigger a warning and fall back to \texttt{book}.

\section{KOMA Class Base}

Each KOMA class carries a distinct set of default options, defined in
\texttt{omnilatex.cls}:

\begin{table}[htbp]
    \centering
    \caption{KOMA class base options}
    \label{tab:koma-base-options}
    \begin{tabular}{@{} l l p{5cm} @{}}
        \toprule
        \textbf{KOMA Class} & \textbf{Options Macro} & \textbf{Default Options} \\
        \midrule
        \texttt{scrbook} &
            \texttt{\textbackslash omnilatex@bookoptions} &
            \texttt{listof=totoc, bibliography=totoc, chapterprefix=true,
            open=right, numbers=noenddot} \\[6pt]
        \texttt{scrreprt} &
            \texttt{\textbackslash omnilatex@reportoptions} &
            \texttt{listof=totoc, bibliography=totoc, chapterprefix=false,
            open=any, numbers=noenddot} \\[6pt]
        \texttt{scrartcl} &
            \texttt{\textbackslash omnilatex@articleoptions} &
            \texttt{bibliography=totoc, numbers=noenddot, parskip=half} \\
        \bottomrule
    \end{tabular}
\end{table}

Two additional option sets exist for specialised doctypes:

\begin{itemize}
    \item \texttt{\textbackslash omnilatex@inlinepaperoptions} -- Same as
          article options with \texttt{titlepage=false}.
    \item \texttt{\textbackslash omnilatex@journaloptions} -- Same as
          article options with \texttt{titlepage=true}.
    \item \texttt{\textbackslash omnilatex@cvoptions} --
          \texttt{numbers=noenddot, parskip=full, fontsize=11pt}.
\end{itemize}

\section{Complete Alias Table}

The following tables list all aliases grouped by their canonical doctype
and KOMA base class.

\subsection{Aliases mapping to \texttt{scrbook}}

\begin{table}[htbp]
    \centering
    \caption{Aliases resolving to \texttt{scrbook}}
    \label{tab:aliases-scrbook}
    \begin{tabular}{@{} l l @{}}
        \toprule
        \textbf{Canonical} & \textbf{Aliases} \\
        \midrule
        book &
            book, books \\[2pt]
        thesis &
            thesis, theses \\[2pt]
        dissertation &
            dissertation, dissertations \\[2pt]
        dictionary &
            dictionary, dictionaries, lexicon, lexicons \\
        \bottomrule
    \end{tabular}
\end{table}

\subsection{Aliases mapping to \texttt{scrreprt}}

\begin{table}[htbp]
    \centering
    \caption{Aliases resolving to \texttt{scrreprt}}
    \label{tab:aliases-scrreprt}
    \begin{tabular}{@{} l l @{}}
        \toprule
        \textbf{Canonical} & \textbf{Aliases} \\
        \midrule
        manual &
            manual, manuals, guide, guides, handbook, handbooks \\[2pt]
        technicalreport &
            report, reports, technicalreport, technical-report,
            technicalreports, technical-reports, techreport,
            tech-report, techreports \\[2pt]
        standard &
            standard, standards \\[2pt]
        patent &
            patent, patents \\[2pt]
        research-proposal &
            research-proposal, researchproposal, research-proposals \\
        \bottomrule
    \end{tabular}
\end{table}

\subsection{Aliases mapping to \texttt{scrartcl}}

\begin{table}[htbp]
    \centering
    \caption{Aliases resolving to \texttt{scrartcl}}
    \label{tab:aliases-scrartcl}
    \begin{tabular}{@{} l l @{}}
        \toprule
        \textbf{Canonical} & \textbf{Aliases} \\
        \midrule
        article &
            article, articles, paper, papers \\[2pt]
        inlinepaper &
            inlinepaper, inlinepapers, inline-research,
            inline-research-paper \\[2pt]
        journal &
            journal, journals, magazine, magazines \\[2pt]
        cv &
            cv, resume, resumes, curriculumvitae \\[2pt]
        cover-letter &
            cover-letter, coverletter \\[2pt]
        poster &
            poster, posters \\[2pt]
        presentation &
            presentation, presentations, slides, talk, talks \\[2pt]
        letter &
            letter, letters \\[2pt]
        homework &
            homework, homeworks \\[2pt]
        exam &
            exam, exams, examination, examinations \\[2pt]
        lecture-notes &
            lecture-notes, lecturenotes, lecture-note \\[2pt]
        syllabus &
            syllabus, syllabi \\[2pt]
        handout &
            handout, handouts \\[2pt]
        memo &
            memo, memos, memorandum, memorandums \\[2pt]
        white-paper &
            white-paper, whitepapers \\[2pt]
        invoice &
            invoice, invoices \\[2pt]
        recipe &
            recipe, recipes \\
        \bottomrule
    \end{tabular}
\end{table}

\section{Canonical Doctypes}

Table~\ref{tab:canonical-doctypes} lists all 26 canonical doctypes with
their KOMA base class and category.  Each canonical doctype has a
corresponding \texttt{.sty} profile in \texttt{config/document-types/}.

\begin{table}[htbp]
    \centering
    \caption{All 26 canonical doctypes}
    \label{tab:canonical-doctypes}
    \begin{tabular}{@{} l l l l @{}}
        \toprule
        \textbf{Canonical} & \textbf{KOMA Base} & \textbf{Category} &
        \textbf{Profile File} \\
        \midrule
        article & scrartcl & Academic & article.sty \\
        inlinepaper & scrartcl & Academic & inlinepaper.sty \\
        journal & scrartcl & Academic & journal.sty \\
        cv & scrartcl & Business & cv.sty \\
        cover-letter & scrartcl & Business & cover-letter.sty \\
        letter & scrartcl & Business & letter.sty \\
        memo & scrartcl & Business & memo.sty \\
        invoice & scrartcl & Business & invoice.sty \\
        white-paper & scrartcl & Business & white-paper.sty \\
        homework & scrartcl & Education & homework.sty \\
        exam & scrartcl & Education & exam.sty \\
        lecture-notes & scrartcl & Education & lecture-notes.sty \\
        syllabus & scrartcl & Education & syllabus.sty \\
        handout & scrartcl & Education & handout.sty \\
        poster & scrartcl & Visual & poster.sty \\
        presentation & scrartcl & Visual & presentation.sty \\
        recipe & scrartcl & Personal & recipe.sty \\
        \midrule
        manual & scrreprt & Reference & manual.sty \\
        technicalreport & scrreprt & Academic & technicalreport.sty \\
        standard & scrreprt & Technical & standard.sty \\
        patent & scrreprt & Technical & patent.sty \\
        research-proposal & scrreprt & Academic & research-proposal.sty \\
        \midrule
        book & scrbook & Publishing & book.sty \\
        thesis & scrbook & Academic & thesis.sty \\
        dissertation & scrbook & Academic & dissertation.sty \\
        dictionary & scrbook & Reference & dictionary.sty \\
        \bottomrule
    \end{tabular}
\end{table}

\section{Creating Custom Doctypes}

New doctypes are added by creating a \texttt{.sty} file in
\texttt{config/document-types/}.  The class automatically loads the
profile whose filename matches the canonical doctype name.

\subsection{Step-by-step}

\begin{enumerate}
    \item Create \texttt{config/document-types/myreport.sty}.

    \item Define the document type metadata:
\begin{minted}{latex}
\NeedsTeXFormat{LaTeX2e}
\ProvidesPackage{config/document-types/myreport}%
    [2026/01/01 OmniLaTeX myreport profile]
\documenttype{My Custom Report}
\end{minted}

    \item Set layout and KOMA options:
\begin{minted}{latex}
\documentfontsize{11pt}
\documentlayout{DIV=12,headinclude=false,footinclude=false}
\KOMAoptions{titlepage=true}
\setcounter{secnumdepth}{3}
\setcounter{tocdepth}{2}
\end{minted}

    \item Configure document appearance:
\begin{minted}{latex}
\documentcolormode{color}
\documentlinespacing{onehalf}
\documentparspacing{half}
\documentfontmode{serif}
\end{minted}

    \item Optionally define a custom \texttt{\textbackslash maketitle}:

\begin{minted}{latex}
\makeatletter
\renewcommand{\maketitle}{%
    \begin{titlepage}
        \centering
        {\huge\bfseries\@title}\par\vspace{1cm}
        {\large\@author}\par\vspace{0.5cm}
        {\normalsize\@date}
    \end{titlepage}
}
\makeatother
\end{minted}

    \item Use with \texttt{\textbackslash documentclass[doctype=myreport]\{omnilatex\}}.
          Optionally add aliases in \texttt{omnilatex.cls} via
          \texttt{\textbackslash omnilatex@matchdoctypes}.
\end{enumerate}

\subsection{Minimal Example}

A minimal custom doctype profile:

\begin{minted}{latex}
\NeedsTeXFormat{LaTeX2e}
\ProvidesPackage{config/document-types/brief}
    [2026/01/01 OmniLaTeX brief profile]

\documenttype{Internal Brief}
\documentfontsize{11pt}
\documentlayout{DIV=12}
\KOMAoptions{titlepage=false}
\setcounter{secnumdepth}{2}
\setcounter{tocdepth}{1}

\documentcolormode{color}
\documentlinespacing{single}
\documentparspacing{half}
\documentfontmode{sans}

\makeatletter
\renewcommand{\maketitle}{%
    \begin{center}
        {\Large\bfseries\@title}\par\medskip
        {\normalsize\@author\quad|\quad\@date}
    \end{center}
    \bigskip\noindent\rule{\textwidth}{0.4pt}\par\medskip
}
\makeatother

\endinput
\end{minted}

\section{Formal Verification}

The doctype-to-class mapping is formally verified in Lean~4.  Two proof
files establish key properties of the system.

\subsection{DoctypeClassMapping.lean}

Located at \texttt{specs/proofs/OmniLaTeXProofs/DoctypeClassMapping.lean},
this file defines three lists of doctype strings grouped by KOMA class
and proves the following theorems:

\begin{description}
    \item[{\texttt{doctype\_class\_deterministic}}] Every doctype maps to
        at most one KOMA class (determinism via \texttt{Option.some.inj}).

    \item[{\texttt{valid\_doctype\_has\_class}}] A representative sample of
        10 doctypes spanning all three classes each resolves to a
        non-\texttt{none} value (proved by \texttt{decide}).

    \item[{\texttt{scrartcl\_count}}] \texttt{scrartclDoctypes.length = 34}.

    \item[{\texttt{scrbook\_count}}] \texttt{scrbookDoctypes.length = 6}.

    \item[{\texttt{scrreprt\_count}}] \texttt{scrreprtDoctypes.length = 20}.

    \item[{\texttt{class\_covering}}] The total list length equals the sum
        of all three class-specific lists
        ($34 + 6 + 20 = 60$).

    \item[{\texttt{scrartcl\_is\_largest}}] \texttt{scrartcl} has strictly
        more doctypes than both \texttt{scrbook} ($34 > 6$) and
        \texttt{scrreprt} ($34 > 20$), proved via \texttt{omega}.
\end{description}

\subsection{DocumentSettings.lean}

Located at \texttt{specs/proofs/OmniLaTeXProofs/DocumentSettings.lean},
this file models the 26 canonical doctypes as an inductive type and
proves structural properties of the class partition:

\begin{description}
    \item[{\texttt{doctype\_class\_deterministic}}] Each canonical doctype
        resolves to exactly one KOMA class (functional correctness).

    \item[{\texttt{class\_partition}}] The 26 doctypes partition into
        article (15), report (6), and book (5), with all three sets
        non-empty: $15 + 6 + 5 = 26$.

    \item[{\texttt{article\_is\_largest}}] The article class contains
        strictly more doctypes than report ($15 > 6$) and book
        ($15 > 5$).

    \item[{\texttt{doctype\_count}}] Total canonical doctype count is 26.
\end{description}
